\documentclass[11pt,a4paper]{article}

\usepackage[utf8]{inputenc}
\usepackage[T1]{fontenc}
\usepackage[polish]{babel}
\usepackage{lmodern}
\usepackage[margin=2.5cm]{geometry}
\usepackage{amsmath,amssymb}
\usepackage{graphicx}
\usepackage{booktabs}
\usepackage{enumitem}
\usepackage{hyperref}
\usepackage{url}

\title{Konspekt projektu badawczego \\
\large Zastosowanie Monte Carlo Tree Search (MCTS/UCT) do gry w warcaby 10$\times$10 bez obowiązkowego bicia z damkami bijącymi po całej diagonali}
\author{Jakub Kindracki \and Wiktor Kobielski}
\date{\today}

\begin{document}
\maketitle

\section{Opis problemu}

Przedmiotem projektu jest zastosowanie algorytmu Monte Carlo Tree Search
(MCTS), w szczególności jego wariantu Upper Confidence Bound applied to
Trees (UCT), do stworzenia sztucznej inteligencji grającej w wybraną grę
dwuosobową z pełną informacją. Grą wybraną do badań są warcaby na planszy
$10\times10$ z modyfikacjami klasycznych zasad: \emph{bicie nie jest
obowiązkowe} oraz \emph{damki (latające) biją wzdłuż całej diagonali}, co
odpowiada zasadom zbliżonym do warcabów międzynarodowych w aspekcie ruchu
damki, lecz różnym w aspekcie obowiązku bicia.

Brak obowiązkowego bicia istotnie zwiększa współczynnik rozgałęzienia
drzewa gry — w każdej pozycji, w której istnieje bicie, gracz może je
wykonać lub wybrać dowolny inny ruch. Oznacza to, że MCTS musi rozważać
znacznie szerszy zbiór akcji niż w klasycznych warcabach, a heurystyki
ukierunkowane na bicia (typowe dla silników warcabowych) mają
ograniczone zastosowanie. Latające damki natomiast wzmacniają wartość
promocji i wprowadzają silne, dalekosiężne interakcje pomiędzy figurami,
co może wpływać na jakość losowych symulacji. Ta kombinacja zasad czyni
grę nietrywialnym poligonem do badania algorytmu MCTS i jego rozszerzeń.

Stan gry reprezentujemy jako planszę $10\times10$ z 20 pionkami po
każdej stronie, ustawionymi na ciemnych polach pierwszych czterech
rzędów każdego z graczy. Gra kończy się zwycięstwem strony, której
przeciwnik nie ma legalnego ruchu (brak figur lub całkowite
zablokowanie). W celu zapewnienia, aby eksperymenty kończyły się w
rozsądnym czasie, wprowadzimy limit liczby półruchów bez bicia i bez
ruchu pionka — przekroczenie limitu skutkuje remisem (zasada zbliżona
do reguły 50 ruchów w szachach).

\section{Wstępny przegląd literatury}

Algorytm MCTS w jego współczesnej postaci został sformalizowany przez
Kocsisa i Szepesváriego~\cite{kocsis2006uct} jako UCT — zastosowanie
strategii UCB1 \cite{auer2002ucb} do równoważenia eksploracji i
eksploatacji w drzewie gry. Browne i in.~\cite{browne2012survey}
przedstawili obszerny przegląd wariantów MCTS i obszarów ich
zastosowań w grach.

Kluczowym osiągnięciem MCTS jest wynik w grze Go: AlphaGo i AlphaZero
\cite{silver2016alphago,silver2017alphazero} pokazały, że MCTS w
połączeniu z głębokimi sieciami neuronowymi prowadzi do nadludzkiego
poziomu gry. W naszym projekcie skupiamy się na "klasycznym" MCTS bez
sieci neuronowych, ze względu na zakres przedmiotu i ograniczenia
sprzętowe.

W ramach usprawnień UCT zbadamy dwa podejścia z literatury:
\begin{itemize}
    \item \textbf{RAVE / AMAF} (Rapid Action Value Estimation, All Moves
        As First) -- Gelly i Silver \cite{gelly2007combining,
        gelly2011monte} zaproponowali heurystykę współdzielenia statystyk
        ruchów: wynik symulacji jest przypisywany nie tylko ruchowi
        faktycznie wykonanemu w węźle, lecz wszystkim ruchom o tym
        samym kluczu, które wystąpiły w dalszej części playoutu.
        Statystyki AMAF są łączone z klasycznymi statystykami UCT za
        pomocą malejącej wagi $\beta(n)$, co znacząco przyspiesza
        zbieżność wartości akcji w rzadko odwiedzanych węzłach.
    \item \textbf{Progressive Bias} -- Chaslot i
        in.~\cite{chaslot2008progressive} wprowadzili wzmocnienie reguły
        wyboru węzła dodatkowym członem heurystycznym
        $w \cdot h(\text{stan})/(n_i+1)$, gdzie $h$ jest statyczną
        funkcją oceny pozycji, a $n_i$ liczbą wizyt dziecka. Wpływ
        członu maleje wraz ze wzrostem $n_i$, dzięki czemu heurystyka
        kieruje przeszukiwaniem we wczesnej fazie, a w miarę gromadzenia
        danych decyduje statystyka Monte~Carlo.
\end{itemize}

Łącznie porównamy zatem trzy warianty MCTS: \textbf{czysty UCT}
(selekcja wg UCB1), \textbf{UCT+RAVE} oraz \textbf{UCT+progressive
bias}. Wszystkie dzielą tę samą strukturę pętli (selekcja --
ekspansja -- losowa symulacja do końca gry -- propagacja wsteczna)
i ten sam mechanizm zrównoleglenia (root parallelization: niezależne
drzewa w procesach roboczych z agregacją liczników wizyt na końcu).

W kontekście warcabów referencyjnym osiągnięciem jest program
\textsc{Chinook} \cite{schaeffer2007checkers}, który dowiódł, że
warcaby angielskie ($8\times8$, z obowiązkowym biciem) są
\emph{rozwiązane} (gra perfekcyjna kończy się remisem). Warcaby
$10\times10$ pozostają znacznie bardziej złożone obliczeniowo i nie
zostały rozwiązane. Programy oparte o $\alpha\beta$-przeszukiwanie z
heurystykami (np. \textsc{Kingsrow}, \textsc{Scan}) dominują w warcabach
międzynarodowych, ale MCTS w tym wariancie gry pozostaje stosunkowo
słabo zbadany w literaturze, co dodatkowo motywuje nasze badania.

Heurystyczna funkcja oceny dla warcabów -- materiał, mobilność, kontrola
centrum, struktura tylnego rzędu -- jest dobrze opisana m.in. w pracach
Samuela~\cite{samuel1959} i pracach autorów Chinook. Wykorzystamy ją
zarówno w graczu czysto heurystycznym (przeciwnik bazowy), jak i w
członie progressive bias.

\paragraph{Gracz heurystyczny.} Bazowy gracz heurystyczny działa w
prostym schemacie \emph{greedy jednopoziomowym}: dla każdego legalnego
ruchu z bieżącej pozycji wykonuje go na kopii stanu, oblicza wartość
statycznej funkcji oceny powstałej pozycji i wybiera ruch maksymalizujący
ocenę (z deterministycznym rozstrzyganiem remisów dzięki ziarnu RNG).
Funkcja oceny jest liniową kombinacją:
\begin{itemize}[noitemsep]
    \item \textbf{materiału} -- pionek o wartości $1{,}0$, damka
        $3{,}0$;
    \item \textbf{zaawansowania pionka} -- mały bonus za bliskość do
        rzędu promocji (premiuje rozwój i potencjał uzyskania damki);
    \item ocena zwracana jest w postaci znormalizowanej
        $\text{biały}/(\text{biały}+\text{czarny}) \in [0,1]$, dzięki
        czemu działa w obu kolorach po prostym odwróceniu znaku.
\end{itemize}
Świadomie zachowujemy heurystykę prostą -- ma być realistycznym, lecz
słabym przeciwnikiem porównawczym i jednocześnie tym samym członem
$h(\text{stan})$ używanym w wariancie progressive bias, co zapewnia
spójność porównań.

\section{Hipotezy badawcze}

Zgodnie z wymaganiami projektu stawiamy cztery hipotezy: trzy
wymagane (porównania graczy) oraz jedną autorską, związaną
bezpośrednio z modyfikacją zasad badanego wariantu warcabów.

\begin{description}[style=nextline,leftmargin=0pt]
    \item[H1 (UCT vs heurystyka).] Czysty UCT z wystarczającym
        budżetem iteracji ($\geq 10\,000$ na ruch) gra istotnie
        skuteczniej od gracza czysto heurystycznego (greedy,
        jednopoziomowy), osiągając $> 60\%$ wygranych w meczu z
        naprzemienną zmianą kolorów. Hipoteza weryfikuje, czy losowe
        symulacje wystarczają do pokonania prostej, ale ukierunkowanej
        wiedzy domenowej w grze o wysokim współczynniku rozgałęzienia.

    \item[H2 (RAVE poprawia UCT).] UCT z rozszerzeniem RAVE wygrywa
        istotnie częściej z czystym UCT przy tym samym budżecie
        iteracji. Motywacja: warcaby 10$\times$10 bez obowiązku bicia
        mają wysoki współczynnik rozgałęzienia, więc współdzielenie
        statystyk ruchów (AMAF) powinno wyraźnie przyspieszyć
        zbieżność oceny rzadko odwiedzanych akcji.

    \item[H3 (Progressive bias poprawia UCT).] UCT z heurystycznym
        progressive bias wygrywa istotnie częściej z czystym UCT,
        a jego przewaga jest największa przy małych budżetach
        iteracji ($\leq 2\,000$/ruch), gdy losowe symulacje nie
        wystarczają do oceny pozycji. Wraz ze wzrostem liczby
        iteracji oczekujemy zaniku przewagi (statystyki MC dominują).

    \item[H4 (hipoteza autorska -- wpływ braku obowiązku bicia).]
        Modyfikacja zasad polegająca na zniesieniu obowiązku bicia
        zwiększa współczynnik rozgałęzienia drzewa gry i tym samym
        \emph{zwiększa względną przewagę} wariantów wspomaganych
        wiedzą (RAVE i progressive bias) nad czystym UCT.
        Konkretnie: zakładamy, że stosunek wygranych \mbox{RAVE
        vs UCT} (i analogicznie progressive bias vs UCT) będzie
        wyższy w naszym wariancie zasad niż w wariancie klasycznym
        (z obowiązkowym biciem) przy tym samym budżecie iteracji.
        Hipoteza dotyka bezpośrednio specyfiki badanej gry i
        odpowiada na pytanie, czy zniesienie heurystycznego
        ,,filtra'' obowiązkowego bicia kompensowane jest skuteczniej
        przez wiedzę dziedzinową, czy przez czysto statystyczne
        wyszukiwanie MC.
\end{description}

\section{Sposób weryfikacji hipotez}

\subsection{Plan eksperymentów}

Każdą hipotezę zweryfikujemy poprzez zorganizowanie turniejów
\emph{round-robin} pomiędzy graczami przy ustalonym budżecie iteracji /
czasu. Dla każdej pary graczy rozegramy $N \geq 50$ partii (po
$N/2$ z każdym kolorem), aby wyeliminować wpływ pierwszego ruchu.

\paragraph{Zbiory danych.} Ponieważ jest to gra dwuosobowa, "danymi"
są wyniki rozegranych partii. Dla zapewnienia powtarzalności:
\begin{itemize}
    \item Każdy gracz przyjmuje parametr \texttt{seed}; ziarna pochodne
        wyprowadzamy deterministycznie z ziarna bazowego.
    \item Każda konfiguracja eksperymentu (para graczy, budżet, wariant
        zasad) zostanie powtórzona z $\geq 5$ różnymi ziarnami bazowymi,
        co pozwoli oszacować rozrzut wyników.
    \item Wszystkie partie zostaną zapisane (sekwencja ruchów + ziarno),
        umożliwiając powtórzenie konkretnego przebiegu.
\end{itemize}

\paragraph{Mierzone wielkości.} Procent wygranych, remisów i przegranych
z 95\%-przedziałem ufności (Wilson lub bootstrap), średnia długość
partii, średni czas na ruch, średnia liczba odwiedzonych węzłów drzewa.

\paragraph{Testy istotności.} Dla porównań dwóch graczy zastosujemy
dwustronny test proporcji (lub bootstrap przedziału ufności na różnicy
procentów wygranych). Hipotezę uznajemy za potwierdzoną, jeśli
$p < 0{,}05$ \emph{i} efekt jest spójny dla różnych ziaren.

\paragraph{Plan opracowania wyników.} Wyniki przedstawimy w postaci:
wykresów słupkowych z przedziałami ufności (porównania
graczy), wykresów krzywych (skuteczność vs budżet iteracji), tabel
zbiorczych w raporcie. Surowe dane zostaną dołączone jako pliki
JSON/CSV.

\subsection{Testy z udziałem ludzi}

Zaimplementujemy interfejs przeglądarkowy (Flask + HTML/JS) umożliwiający
grę człowiek-komputer. Przeprowadzimy testy z co najmniej 5 osobami
(różny poziom znajomości warcabów). Każda osoba rozegra po jednej
partii z każdym wariantem AI; zbierzemy ankietę dotyczącą subiektywnej
oceny "ludzkiego" charakteru i siły gry, a także wyniki obiektywne.

\section{Harmonogram działań}

\begin{tabular}{p{3.2cm}p{11cm}}
\toprule
\textbf{Termin} & \textbf{Zadanie} \\
\midrule
Tydzień 1--2 & Implementacja silnika gry (warcaby 10$\times$10 z
                wybranymi zasadami), testy jednostkowe generatora ruchów. \\
Tydzień 3 & Implementacja czystego UCT i interfejsu gracza, gracza losowego
            i heurystycznego. \\
Tydzień 4 & Implementacja wariantów UCT: RAVE i progressive bias. \\
Tydzień 5 & Implementacja interfejsu przeglądarkowego, runner
            eksperymentalny CLI z obsługą ziarna i raportowania. \\
Tydzień 6 & Eksperymenty H1--H3 (porównania graczy AI). Optymalizacja
            wydajności (parallel root MCTS). \\
Tydzień 7 & Eksperymenty H4 (porównanie wariantów zasad), testy z
            udziałem 5+ osób. \\
Tydzień 8 & Opracowanie wyników, pisanie raportu, przygotowanie
            prezentacji. \\
\bottomrule
\end{tabular}

\section{Ogólny projekt techniczny}

Całość projektu zostanie napisana w języku \textbf{Python~3}
(cross-platform, jeden ekosystem dla silnika gry, MCTS, runnera
eksperymentów oraz interfejsu webowego). Korzystamy wyłącznie ze
standardowej biblioteki oraz kilku popularnych pakietów:

\begin{itemize}[noitemsep]
    \item \texttt{multiprocessing} (stdlib) -- root parallelization
        MCTS (niezależne drzewa w procesach roboczych, agregacja
        liczników wizyt na końcu).
    \item \texttt{Flask} -- lekki serwer webowy obsługujący interfejs
        gry człowiek--komputer (REST + HTML/JS).
    \item \texttt{matplotlib} / \texttt{pandas} -- analiza wyników i
        wykresy do raportu końcowego.
\end{itemize}

\paragraph{Reprodukowalność.} Wszyscy gracze przyjmują \texttt{seed}.
Eksperymenty alternują kolory i używają ziaren wyprowadzonych
deterministycznie z bazowego. Pełna konfiguracja partii (gracze, ziarno,
zasady) zapisywana jest w wynikach, co pozwala odtworzyć dowolny
przebieg.

\paragraph{Wykorzystanie LLM.} Duże modele językowe (Claude Code) były i
będą wykorzystywane jako asystent programistyczny -- generowanie
szkieletu kodu, refaktoryzacja, sugerowanie testów. Każde użycie zostanie
udokumentowane i podsumowane w raporcie końcowym, zgodnie z wymaganiem
projektu. Ostateczna treść raportu, hipotezy, projekt eksperymentów i
analiza wyników są dziełem autorów.

\begin{thebibliography}{9}

\bibitem{kocsis2006uct}
L.~Kocsis and C.~Szepesvári,
\emph{Bandit based Monte-Carlo Planning},
ECML, 2006.

\bibitem{auer2002ucb}
P.~Auer, N.~Cesa-Bianchi, P.~Fischer,
\emph{Finite-time Analysis of the Multiarmed Bandit Problem},
Machine Learning, 47(2--3):235--256, 2002.

\bibitem{browne2012survey}
C.~Browne et al.,
\emph{A Survey of Monte Carlo Tree Search Methods},
IEEE Trans. on Computational Intelligence and AI in Games, 4(1):1--43,
2012.

\bibitem{gelly2007combining}
S.~Gelly and D.~Silver,
\emph{Combining online and offline knowledge in UCT},
ICML, 2007.

\bibitem{gelly2011monte}
S.~Gelly and D.~Silver,
\emph{Monte-Carlo tree search and rapid action value estimation in
computer Go},
Artificial Intelligence, 175(11):1856--1875, 2011.

\bibitem{chaslot2008progressive}
G.~M.~J.-B.~Chaslot, M.~H.~M.~Winands, H.~J.~van~den~Herik,
J.~W.~H.~M.~Uiterwijk, B.~Bouzy,
\emph{Progressive strategies for Monte-Carlo tree search},
New Mathematics and Natural Computation, 4(3):343--357, 2008.

\bibitem{silver2016alphago}
D.~Silver et al.,
\emph{Mastering the game of Go with deep neural networks and tree
search},
Nature, 529(7587):484--489, 2016.

\bibitem{silver2017alphazero}
D.~Silver et al.,
\emph{Mastering Chess and Shogi by Self-Play with a General
Reinforcement Learning Algorithm},
arXiv:1712.01815, 2017.

\bibitem{schaeffer2007checkers}
J.~Schaeffer, N.~Burch, Y.~Björnsson, A.~Kishimoto, M.~Müller,
R.~Lake, P.~Lu, S.~Sutphen,
\emph{Checkers is Solved},
Science, 317(5844):1518--1522, 2007.

\bibitem{samuel1959}
A.~L.~Samuel,
\emph{Some Studies in Machine Learning Using the Game of Checkers},
IBM Journal of Research and Development, 3(3):210--229, 1959.

\end{thebibliography}

\end{document}
